%% ============================================================================
%%  back/cv.tex
%%  کارنامک (رزومه علمی و تحصیلی نویسنده)
%%  طراحی شده برای قالب رسمی دانشگاه خوارزمی
%% ============================================================================
%%
%%  توضیح:
%%  این صفحه برای معرفی سوابق تحصیلی، شغلی، زمینه‌های پژوهشی و افتخارات علمی
%%  دانش‌آموخته در پایان پایان‌نامه اختصاص یافته است.
%% ============================================================================

\makeCV{%
\noindent
\textbf{۱. مشخصات فردی:}\par
\begin{itemize}[leftmargin=1.5cm, noitemsep]
  \item \textbf{نام و نام خانوادگی:} آیریک پژوهشگر
  \item \textbf{سال تولد:} ۱۳۷۸
  \item \textbf{محل تولد:} تهران
  \item \textbf{رایانامه:} \lr{\texttt{researcher@khu.ac.ir}}
\end{itemize}

\vspace{0.4cm}
\noindent
\textbf{۲. سوابق تحصیلی:}\par
\begin{itemize}[leftmargin=1.5cm, noitemsep]
  \item \textbf{کارشناسی ارشد:} مهندسی کامپیوتر - هوش مصنوعی، دانشگاه خوارزمی، تهران (۱۴۰۱--۱۴۰۳)\par
  \textit{عنوان پایان‌نامه:} طراحی و پیاده‌سازی قالب جامع پایان‌نامه دانشگاه خوارزمی در سامانه لاتک
  \item \textbf{کارشناسی:} مهندسی کامپیوتر، دانشگاه خوارزمی، تهران (۱۳۹۷--۱۴۰۱)
  \item \textbf{دیپلم:} علوم ریاضی و فیزیک، دبیرستان نمونه دولتی، تهران (۱۳۹۳--۱۳۹۷)
\end{itemize}

\vspace{0.4cm}
\noindent
\textbf{۳. زمینه‌های پژوهشی و تخصصی:}\par
\begin{itemize}[leftmargin=1.5cm, noitemsep]
  \item حروف‌چینی علمی رایانه‌ای و سامانه‌های متن‌باز (\lr{LaTeX, XePersian})
  \item یادگیری ماشین و شبکه‌های عصبی عمیق (\lr{Machine Learning \& Deep Learning})
  \item پردازش زبان‌های طبیعی و مدل‌های زبانی (\lr{NLP \& LLMs})
  \item بهینه‌سازی محاسباتی و الگوریتم‌های گرادیانی
\end{itemize}

\vspace{0.4cm}
\noindent
\textbf{۴. افتخارات و جوایز علمی:}\par
\begin{itemize}[leftmargin=1.5cm, noitemsep]
  \item رتبه برتر علمی در مقطع کارشناسی ارشد مهندسی کامپیوتر دانشگاه خوارزمی (۱۴۰۳)
  \item پذیرفته‌شده بدون آزمون استعدادهای درخشان در مقطع دکتری تخصصی (۱۴۰۳)
  \item رتبه برتر در مسابقات برنامه‌نویسی دانشگاهی (۱۴۰۰)
\end{itemize}

\vspace{0.4cm}
\noindent
\textbf{۵. مهارت‌های فنی و نرم‌افزاری:}\par
\begin{itemize}[leftmargin=1.5cm, noitemsep]
  \item \textbf{زبان‌های برنامه‌نویسی:} \lr{Python, C++, Java, MATLAB, SQL}
  \item \textbf{ابزارهای یادگیری عمیق:} \lr{PyTorch, TensorFlow, Hugging Face}
  \item \textbf{ابزارهای حروف‌چینی و توسعه:} \lr{XeLaTeX, Git, Linux, Docker}
\end{itemize}
}
